% ============================================================
%  VisiSize — Documentación Técnica Detallada
%  Flujograma y descripción de procesos
%  Fecha: Marzo 2026
% ============================================================
\documentclass[12pt, a4paper]{article}

% ── Codificación y lenguaje ─────────────────────────────────
\usepackage[utf8]{inputenc}
\usepackage[T1]{fontenc}
\usepackage[spanish]{babel}

% ── Tipografía y espaciado ──────────────────────────────────
\usepackage{lmodern}
\usepackage{microtype}
\usepackage{setspace}
\onehalfspacing

% ── Márgenes ────────────────────────────────────────────────
\usepackage[top=2.5cm, bottom=2.5cm, left=3cm, right=2.5cm]{geometry}

% ── Gráficos y colores ──────────────────────────────────────
\usepackage{tikz}
\usetikzlibrary{shapes.geometric, arrows.meta, positioning, fit,
                backgrounds, shadows, calc, decorations.pathreplacing}
\usepackage{xcolor}
\usepackage{graphicx}

% ── Tablas ──────────────────────────────────────────────────
\usepackage{booktabs}
\usepackage{tabularx}
\usepackage{longtable}
\usepackage{array}
\usepackage{multirow}

% ── Matemáticas ─────────────────────────────────────────────
\usepackage{amsmath}
\usepackage{amssymb}

% ── Código fuente ────────────────────────────────────────────
\usepackage{listings}
\usepackage{courier}

% ── Citaciones IEEE ──────────────────────────────────────────
\usepackage{cite}
\usepackage{url}

% ── Cajas y destacados ───────────────────────────────────────
\usepackage{tcolorbox}
\tcbuselibrary{skins, breakable, listingsutf8}

% ── Encabezados y pies de página ─────────────────────────────
\usepackage{fancyhdr}
\pagestyle{fancy}
\fancyhf{}
\fancyhead[L]{\small VisiSize — Documentación Técnica}
\fancyhead[R]{\small Marzo 2026}
\fancyfoot[C]{\thepage}
\renewcommand{\headrulewidth}{0.4pt}

% ── Hipervínculos ────────────────────────────────────────────
\usepackage{hyperref}
\hypersetup{
    colorlinks=true,
    linkcolor=blue!60!black,
    urlcolor=blue!70!black,
    citecolor=green!50!black,
    pdfauthor={JRapp — UDC 2026-1},
    pdftitle={VisiSize: Documentación Técnica y Flujogramas},
    pdfsubject={Medición de área por visión artificial},
}

% ── Títulos de sección ───────────────────────────────────────
\usepackage{titlesec}
\titleformat{\section}[block]
  {\Large\bfseries\color{blue!60!black}}{\thesection.}{0.8em}{}[\vspace{-0.3em}\hrule\vspace{0.3em}]
\titleformat{\subsection}[block]
  {\large\bfseries\color{blue!40!black}}{\thesubsection.}{0.6em}{}
\titleformat{\subsubsection}[block]
  {\normalsize\bfseries\color{gray!70!black}}{\thesubsubsection.}{0.5em}{}

% ── Paleta de colores del proyecto ───────────────────────────
\definecolor{visiblu}{RGB}{25, 90, 170}
\definecolor{visigray}{RGB}{80, 90, 100}
\definecolor{catA}{RGB}{76, 175, 80}      % verde
\definecolor{catB}{RGB}{33, 150, 243}     % azul
\definecolor{catC}{RGB}{255, 152, 0}      % ámbar
\definecolor{catD}{RGB}{255, 87, 34}      % naranja
\definecolor{catErr}{RGB}{244, 67, 54}    % rojo
\definecolor{pipebg}{RGB}{240, 248, 255}
\definecolor{pipeborder}{RGB}{25, 90, 170}
\definecolor{notecolor}{RGB}{255, 248, 220}
\definecolor{noteborder}{RGB}{204, 153, 0}
\definecolor{codebg}{RGB}{245, 245, 245}
\definecolor{catE}{RGB}{156, 39, 176}     % morado — Zona E
\definecolor{catOut}{RGB}{244, 67, 54}    % rojo   — Fuera de rango

% ── Caja informativa ─────────────────────────────────────────
\newtcolorbox{infobox}[1][]{
  enhanced, breakable,
  colback=pipebg, colframe=pipeborder,
  fonttitle=\bfseries\color{white},
  coltitle=white, attach boxed title to top left={yshift=-2mm, xshift=4mm},
  boxed title style={colback=visiblu, rounded corners},
  rounded corners, #1
}

\newtcolorbox{notebox}[1][]{
  enhanced, breakable,
  colback=notecolor, colframe=noteborder,
  fonttitle=\bfseries, #1
}

\newtcolorbox{formulabox}[1][]{
  enhanced,
  colback=blue!4!white, colframe=visiblu,
  fonttitle=\bfseries\color{visiblu},
  attach boxed title to top left={yshift=-2mm, xshift=5mm},
  boxed title style={colback=blue!15!white, colframe=visiblu},
  rounded corners, #1
}

% ── Estilo de código ─────────────────────────────────────────
\lstdefinestyle{pycode}{
  backgroundcolor=\color{codebg},
  basicstyle=\footnotesize\ttfamily,
  keywordstyle=\color{blue!70!black}\bfseries,
  commentstyle=\color{green!50!black}\itshape,
  stringstyle=\color{red!60!black},
  breaklines=true,
  frame=single,
  rulecolor=\color{gray!40},
  numbers=left,
  numberstyle=\tiny\color{gray},
  stepnumber=1,
  language=Python,
  showstringspaces=false,
}

% ── TikZ: estilos de cajas para flujograma ──────────────────
\tikzset{
  startstop/.style={
    rectangle, rounded corners=8pt, minimum width=3.5cm, minimum height=0.9cm,
    draw=visiblu!80, fill=visiblu!20, text centered, font=\small\bfseries,
    drop shadow
  },
  process/.style={
    rectangle, minimum width=4.5cm, minimum height=0.85cm,
    draw=blue!50!black, fill=blue!8, text centered, font=\small,
    drop shadow
  },
  decision/.style={
    diamond, aspect=2.5, minimum width=4cm, minimum height=0.8cm,
    draw=orange!80!black, fill=orange!15, text centered, font=\small,
    inner sep=1pt, drop shadow
  },
  io/.style={
    trapezium, trapezium left angle=70, trapezium right angle=110,
    minimum width=3.8cm, minimum height=0.8cm,
    draw=green!60!black, fill=green!10, text centered, font=\small,
    drop shadow
  },
  storage/.style={
    cylinder, shape border rotate=90, aspect=0.15,
    minimum height=1cm, minimum width=2.8cm,
    draw=purple!60!black, fill=purple!10, text centered, font=\small,
    drop shadow
  },
  arrow/.style={-Stealth, thick, color=visigray},
  warrow/.style={-Stealth, thick, color=orange!70!black, dashed},
}

% ============================================================
\begin{document}
% ============================================================

% ── Portada ─────────────────────────────────────────────────
\begin{titlepage}
  \centering
  \vspace*{1.5cm}

  {\Huge\bfseries\color{visiblu} VisiSize}\\[0.4cm]
  {\large\color{visigray} Sistema de Medición de Área por Visión Artificial}\\[1.5cm]

  \begin{tikzpicture}
    \draw[visiblu, line width=3pt, rounded corners=12pt]
      (0,0) rectangle (12,3);
    \fill[visiblu!10, rounded corners=12pt]
      (0.1,0.1) rectangle (11.9,2.9);
    \node[text width=11.5cm, align=center] at (6,1.5) {%
      \large\color{visigray}%
      Documentación técnica detallada de arquitectura,\\
      flujogramas y algoritmos del sistema de clasificación\\
      de retales metálicos en \textbf{ServiMetales de Colombia S.A.S.}
    };
  \end{tikzpicture}

  \vspace{1.5cm}
  {\large\bfseries Documentación Técnica y Pedagógica}\\[0.5cm]
  {\normalsize Versión 1.0.0}\\[2cm]

  \begin{tabular}{rl}
    \textbf{Cliente:}    & ServiMetales de Colombia S.A.S. \\[0.3cm]
    \textbf{Proyecto:}   & JRapp --- AreaCamPython\_v1 \\[0.3cm]
    \textbf{Versión:}    & Prototipo v1.0 \\[0.3cm]
    \textbf{Semestre:}   & UDC 2026-1 \\[0.3cm]
    \textbf{Tecnologías:} & Python 3 $\cdot$ OpenCV $\cdot$ SQLite $\cdot$ CustomTkinter \\[0.3cm]
    \textbf{Fecha:}      & Marzo 2026 \\
  \end{tabular}

  \vfill
  {\small\color{gray} Este documento describe el funcionamiento interno
  del sistema VisiSize de manera detallada, incluyendo flujogramas,\\
  fórmulas matemáticas y explicaciones paso a paso de cada proceso.}
\end{titlepage}

% ── Tabla de contenidos ──────────────────────────────────────
\tableofcontents
\newpage

% ============================================================
\section{Introducción: ¿Qué es VisiSize?}
% ============================================================

\begin{infobox}[title={Concepto Central --- ServiMetales de Colombia S.A.S.}]
\textbf{VisiSize} (prototipo v1.0) es una aplicación de escritorio que utiliza
visión artificial para \textbf{medir el área real} de retales metálicos
resultantes del proceso de corte y manufactura en \textbf{ServiMetales de Colombia S.A.S.},
expresando dicha área en metros cuadrados ($\text{m}^2$) y asignando automáticamente
una \textbf{zona de almacenamiento} dentro de la planta industrial, todo sin contacto
físico con la pieza \cite{malamas2003}.
\end{infobox}

\vspace{0.5cm}

En el sector metalmecánico, los procesos de corte de láminas, perfiles y tubos de acero
generan \textbf{retales metálicos} de diversas dimensiones que deben gestionarse
eficientemente para su reutilización o disposición final. VisiSize automatiza la
inspección visual de estos retales: la cámara apunta a la pieza, el algoritmo detecta
su contorno, calcula su área y la asigna digitalmente a la zona correspondiente \cite{chin1982}.

\subsection{¿Cómo lo logra? --- La idea básica}

El proceso tiene tres ingredientes fundamentales:

\begin{enumerate}
  \item \textbf{Ver:} La cámara captura imágenes en tiempo real (30 fotogramas por segundo).
  \item \textbf{Detectar:} Un algoritmo de procesamiento de imagen identifica el contorno del objeto.
  \item \textbf{Medir:} Usando un factor de conversión previamente calculado (la \emph{calibración}),
        convierte el área en píxeles a área en metros cuadrados.
\end{enumerate}

\begin{notebox}[title={Nota}]
Un \textbf{píxel} es el punto más pequeño que forma una imagen digital.
Si tienes una imagen de 1280×720 píxeles, tiene 921\,600 píxeles en total.
El reto del sistema es saber cuántos metros reales corresponde a cada píxel.
\end{notebox}

\subsection{Arquitectura general del sistema}

El proyecto está organizado en capas bien definidas (arquitectura limpia):

\vspace{0.5cm}
\begin{center}
\begin{tikzpicture}[node distance=0.4cm]
  \node[draw=visiblu, fill=visiblu!10, rounded corners=6pt,
        minimum width=10cm, minimum height=0.7cm, font=\small\bfseries]
        (ui) {\texttt{ui/} — Interfaz de usuario (CustomTkinter)};
  \node[draw=blue!60, fill=blue!8, rounded corners=6pt,
        minimum width=10cm, minimum height=0.7cm, font=\small\bfseries,
        below=of ui]
        (svc) {\texttt{services/} — Lógica de negocio y orquestación};
  \node[draw=green!50!black, fill=green!8, rounded corners=6pt,
        minimum width=10cm, minimum height=0.7cm, font=\small\bfseries,
        below=of svc]
        (core) {\texttt{core/} — Fórmulas, entidades y constantes puras};
  \node[draw=purple!60, fill=purple!8, rounded corners=6pt,
        minimum width=10cm, minimum height=0.7cm, font=\small\bfseries,
        below=of core]
        (data) {\texttt{data/} — Base de datos SQLite y archivos JSON};

  \draw[arrow] (ui) -- (svc);
  \draw[arrow] (svc) -- (core);
  \draw[arrow] (svc) -- (data);

  \node[right=0.5cm of ui, font=\scriptsize\color{gray}] {Ventanas, widgets, tema visual};
  \node[right=0.5cm of svc, font=\scriptsize\color{gray}] {Cámara, calibración, medición, OpenCV};
  \node[right=0.5cm of core, font=\scriptsize\color{gray}] {FactorK, categorías, validaciones};
  \node[right=0.5cm of data, font=\scriptsize\color{gray}] {Mediciones, calibración, imágenes};
\end{tikzpicture}
\end{center}

\vspace{0.3cm}
La dirección de las flechas indica la dependencia: la UI depende de los servicios,
los servicios dependen del núcleo y de los datos, pero \textbf{el núcleo no sabe nada
de la UI ni de la base de datos} --- esto hace el código más mantenible y testeable.

\subsection{Contexto industrial: ServiMetales de Colombia S.A.S.}

ServiMetales de Colombia S.A.S. es una empresa del sector metalmecánico dedicada al
procesamiento, corte y distribución de aceros. Sus procesos productivos generan
\textbf{retales metálicos} --- fragmentos de láminas, tubos, perfiles estructurales y placas
--- que deben gestionarse sistemáticamente para maximizar su reutilización y minimizar
los residuos industriales \cite{malamas2003, chin1982}.

Los objetivos del sistema en esta versión prototipo son:

\begin{itemize}
  \item \textbf{Medir} autonomáticamente el área de retales con visión artificial \cite{gonzalez2018}.
  \item \textbf{Clasificar} cada pieza en una de cinco zonas de almacenamiento (A--E)
        según su área calculada.
  \item \textbf{Inventariar} el material potencialmente reutilizable en una base de datos SQLite \cite{owens2010}.
  \item \textbf{Sentar las bases} para un gemelo digital del layout de planta en versiones futuras \cite{tao2018}.
\end{itemize}

\begin{infobox}[title={Definición operacional: Fuera de Rango (FdR)}]
En el prototipo v1.0, una pieza se considera \textbf{fuera de rango} en dos situaciones:
\begin{itemize}
  \item \textbf{Demasiado peque\~na:} área real $< 0{,}0025\,\text{m}^2$
        (equivalente a $5\,\text{cm} \times 5\,\text{cm}$). La pieza no es
        procesable o almacenable de forma individual.
  \item \textbf{Demasiado grande:} la pieza excede el \textbf{campo de visión
        completo de la cámara} (su proyección ocupa más del 95\,\% del fotograma).
        Su contorno no puede delimitarse con precisión.
\end{itemize}
Las piezas FdR se se\~nalizan en \textcolor{catOut}{\textbf{rojo}} en la interfaz
y \textbf{no se asignan a ninguna zona de almacenamiento}.
\end{infobox}

\subsubsection*{Zonas de almacenamiento --- Prototipo v1.0}

El sistema define cinco zonas ordenadas alfabéticamente. En esta versión, la
asignación se basa exclusivamente en el área calculada. A futuro, el área
será un parámetro más junto con el tipo de pieza detectado por el clasificador
YOLO (lámina, tubo, perfil, placa) \cite{li2020, tao2018}.

\vspace{0.3cm}
\begin{center}
\begin{tabular}{p{1.1cm} p{2.8cm} p{3.8cm} p{2.5cm} p{1.8cm}}
\toprule
\textbf{Zona} & \textbf{Descripción} & \textbf{Rango de área (v1.0)} & \textbf{Color digital} & \textbf{Código hex} \\
\midrule
\textbf{A} & Primera zona & $[0{,}0025 - 1{,}0)\,\text{m}^2$ &
  \textcolor{catA}{\rule{1.2cm}{0.25cm}} Verde & \texttt{\#4caf50} \\[3pt]
\textbf{B} & Segunda zona & $[1{,}0 - 2{,}0)\,\text{m}^2$ &
  \textcolor{catB}{\rule{1.2cm}{0.25cm}} Azul & \texttt{\#2196f3} \\[3pt]
\textbf{C} & Tercera zona & $[2{,}0 - 3{,}0)\,\text{m}^2$ &
  \textcolor{catC}{\rule{1.2cm}{0.25cm}} Ámbar & \texttt{\#ff9800} \\[3pt]
\textbf{D} & Cuarta zona & $[3{,}0 - 4{,}0)\,\text{m}^2$ &
  \textcolor{catD}{\rule{1.2cm}{0.25cm}} Naranja & \texttt{\#ff5722} \\[3pt]
\textbf{E} & Quinta zona & $\geq 4{,}0\,\text{m}^2$ (cabe en frame) &
  \textcolor{catE}{\rule{1.2cm}{0.25cm}} Morado & \texttt{\#9c27b0} \\
\midrule
\textbf{FdR} & Fuera de rango & $< 0{,}0025\,\text{m}^2$ \textit{ó} excede FOV &
  \textcolor{catOut}{\rule{1.2cm}{0.25cm}} Rojo & \texttt{\#f44336} \\
\bottomrule
\end{tabular}
\end{center}

\begin{notebox}[title={Evolución futura de la clasificación}]
En versiones posteriores, el módulo YOLO identificará el tipo de pieza metálica
(lámina, tubo, perfil L/C/H, placa). La asignación de zona incorporará el tipo,
la geometría (circularidad, relación de aspecto) y el área como parámetros de
decisión conjuntos, habilitando el gemelo digital del layout \cite{tao2018, jocher2023}.
\end{notebox}

% ============================================================
\section{Flujograma General del Sistema}
% ============================================================

El siguiente diagrama muestra el flujo completo de vida de una medición,
desde que el usuario abre la aplicación hasta que el dato queda guardado.

\vspace{0.8cm}
\begin{center}
\begin{tikzpicture}[node distance=0.6cm and 1.5cm]

  % Nodos
  \node[startstop] (inicio) {INICIO};

  \node[process, below=of inicio]
    (config) {Cargar configuración\\(\texttt{config.py})};

  \node[process, below=of config]
    (calibcheck) {Verificar calibración\\(\texttt{calibration.json})};

  \node[decision, below=of calibcheck]
    (calexists) {¿Calibración\\válida?};

  \node[io, right=2.5cm of calexists]
    (calwiz) {Asistente de\\calibración (3 pasos)};

  \node[process, below=of calexists]
    (camstart) {Inicializar cámara\\(\texttt{CameraService})};

  \node[decision, below=of camstart]
    (camok) {¿Cámara\\disponible?};

  \node[io, right=2.5cm of camok]
    (camerr) {Mostrar error\\y retry};

  \node[process, below=of camok]
    (pipeline) {Ejecutar pipeline de\\imagen por frame\\(\texttt{OpenCVService})};

  \node[decision, below=of pipeline]
    (objdet) {¿Objeto\\detectado?};

  \node[io, left=2.5cm of objdet]
    (noobj) {Mostrar mensaje\\«sin objeto»};

  \node[process, below=of objdet]
    (calcarea) {Calcular área real\\con FactorK corregido};

  \node[process, below=of calcarea]
    (classify) {Verificar rango y asignar\\zona A/B/C/D/E o FdR};

  \node[process, below=of classify]
    (overlay) {Mostrar resultado en\\tiempo real (overlay)};

  \node[decision, below=of overlay]
    (capture) {¿Usuario presiona\\«Capturar»?};

  \node[process, below=of capture]
    (save) {Guardar medición\\en SQLite + PNG};

  \node[startstop, below=of save]
    (fin) {FIN del ciclo\\(vuelve al loop)};

  % Flechas principales
  \draw[arrow] (inicio) -- (config);
  \draw[arrow] (config) -- (calibcheck);
  \draw[arrow] (calibcheck) -- (calexists);
  \draw[arrow] (calexists) -- node[left, font=\scriptsize]{Sí} (camstart);
  \draw[arrow] (calexists) -- node[above, font=\scriptsize]{No} (calwiz);
  \draw[arrow] (calwiz) |- (camstart);
  \draw[arrow] (camstart) -- (camok);
  \draw[arrow] (camok) -- node[left, font=\scriptsize]{Sí} (pipeline);
  \draw[arrow] (camok) -- node[above, font=\scriptsize]{No} (camerr);
  \draw[arrow] (camerr) |- (camstart);
  \draw[arrow] (pipeline) -- (objdet);
  \draw[arrow] (objdet) -- node[left, font=\scriptsize]{Sí} (calcarea);
  \draw[arrow] (objdet) -- node[above, font=\scriptsize]{No} (noobj);
  \draw[arrow] (noobj) |- (pipeline);
  \draw[arrow] (calcarea) -- (classify);
  \draw[arrow] (classify) -- (overlay);
  \draw[arrow] (overlay) -- (capture);
  \draw[arrow] (capture) -- node[left, font=\scriptsize]{Sí} (save);
  \draw[arrow] (save) -- (fin);
  \draw[warrow] (capture.east) -- ++(1.8,0) |- node[right, font=\scriptsize]{No} (pipeline.east);

\end{tikzpicture}
\end{center}

% ============================================================
\section{Proceso 1: Inicialización del Sistema}
% ============================================================

Cuando el usuario ejecuta \texttt{main.py}, la aplicación realiza una serie de
pasos de preparación antes de mostrar la interfaz:

\vspace{0.5cm}
\begin{center}
\begin{tikzpicture}[node distance=0.5cm]
  \node[startstop] (s) {\texttt{main.py} ejecutado};
  \node[process, below=of s] (log) {Configurar sistema de logging\\(\texttt{app.log})};
  \node[process, below=of log] (dirs) {Crear directorios si no existen\\(\texttt{data/}, \texttt{data/images/})};
  \node[process, below=of dirs] (db) {Inicializar base de datos SQLite\\(\texttt{measurements.db})};
  \node[process, below=of db] (ui) {Instanciar ventana principal\\(\texttt{MainWindow})};
  \node[process, below=of ui] (loop) {Iniciar bucle de eventos Tkinter};

  \draw[arrow] (s)--(log)--(dirs)--(db)--(ui)--(loop);
\end{tikzpicture}
\end{center}

\vspace{0.4cm}
\begin{infobox}[title=¿Qué es el sistema de logging?]
El \textbf{logging} es como un diario de actividades del programa. Cada vez que algo
importante sucede (un error, una calibración exitosa, una medición guardada), queda
registrado con fecha y hora en el archivo \texttt{app.log}. Esto permite diagnosticar
problemas sin necesidad de tener el programa abierto.
\end{infobox}

% ============================================================
\section{Proceso 2: Calibración — El Puente entre Píxeles y Metros}
% ============================================================

La calibración es el proceso más crítico del sistema.
\textbf{Sin calibración no hay medición válida.}

\subsection{¿Por qué es necesaria?}

Piensa en esto: si pones una hoja A4 a 50 cm de la cámara, ocupa, digamos, 120.000 píxeles.
Pero si la alejas a 100 cm, esa \textit{misma} hoja ahora solo ocupa 30.000 píxeles.
El área real no cambió, pero la imagen sí. La calibración resuelve este problema al establecer
una relación matemática entre píxeles y metros para una distancia específica.

\subsection{El Factor de Conversión (FactorK)}

\begin{formulabox}[title=Fórmula del FactorK]
\vspace{0.3cm}
$$\text{FactorK} = \frac{\text{Área real conocida}\ (\text{m}^2)}{\text{Área medida en píxeles}\ (\text{px}^2)}$$
\vspace{0.1cm}

\textbf{Ejemplo numérico:} Una hoja A4 tiene área real de $0{,}0623\ \text{m}^2$.
Si en la imagen ocupa 48.500 px²:
$$\text{FactorK} = \frac{0{,}0623}{48500} \approx 1{,}285 \times 10^{-6}\ \frac{\text{m}^2}{\text{px}^2}$$
\end{formulabox}

\subsection{Flujograma del asistente de calibración (3 pasos)}

\vspace{0.5cm}
\begin{center}
\begin{tikzpicture}[node distance=0.55cm and 2cm]

  \node[startstop] (s) {Iniciar calibración};

  % PASO 1
  \node[process, below=0.7cm of s]
    (p1a) {\textbf{Paso 1:} Capturar frame actual};
  \node[process, below=of p1a]
    (p1b) {Congelar imagen de la cámara};
  \node[io, below=of p1b]
    (p1c) {Mostrar imagen congelada al usuario};

  % PASO 2
  \node[process, below=0.7cm of p1c]
    (p2a) {\textbf{Paso 2:} Seleccionar objeto de referencia};
  \node[io, below=of p2a]
    (p2b) {Usuario dibuja rectángulo\\con el ratón};
  \node[process, below=of p2b]
    (p2c) {Calcular área del rectángulo en px²\\$A_{px} = \text{ancho} \times \text{alto}$};
  \node[decision, below=of p2c]
    (p2d) {$A_{px} \geq 1000\ \text{px}^2$?};
  \node[io, right=2cm of p2d]
    (p2err) {Advertencia:\\área muy pequeña};

  % PASO 3
  \node[process, below=0.7cm of p2d]
    (p3a) {\textbf{Paso 3:} Ingresar datos reales};
  \node[io, below=of p3a]
    (p3b) {Usuario escribe:\\área real (m²) + distancia (cm)};
  \node[process, below=of p3b]
    (p3c) {Calcular FactorK base};
  \node[process, below=of p3c]
    (p3d) {Validar rango del FactorK\\$[10^{-9},\ 1000]$};
  \node[decision, below=of p3d]
    (p3val) {¿FactorK\\válido?};
  \node[io, right=2cm of p3val]
    (p3err) {Error de validación};
  \node[storage, below=0.7cm of p3val]
    (save) {Guardar en\\calibration.json};
  \node[startstop, below=of save]
    (fin) {Calibración completada};

  % Arrows
  \draw[arrow] (s)--(p1a)--(p1b)--(p1c);
  \draw[arrow] (p1c)--(p2a)--(p2b)--(p2c)--(p2d);
  \draw[arrow] (p2d) -- node[left, font=\scriptsize]{Sí} (p3a);
  \draw[arrow] (p2d) -- node[above, font=\scriptsize]{No} (p2err);
  \draw[arrow] (p2err) |- (p2b);
  \draw[arrow] (p3a)--(p3b)--(p3c)--(p3d)--(p3val);
  \draw[arrow] (p3val) -- node[left, font=\scriptsize]{Sí} (save);
  \draw[arrow] (p3val) -- node[above, font=\scriptsize]{No} (p3err);
  \draw[arrow] (p3err) |- (p3b);
  \draw[arrow] (save)--(fin);

\end{tikzpicture}
\end{center}

\subsection{Corrección por distancia}

Si durante la medición el usuario ajusta la distancia cámara-objeto
(por ejemplo, acerca la cámara), el FactorK se recalcula en tiempo real:

\begin{formulabox}[title=Corrección cuadrática por distancia]
$$\text{FactorK}_{\text{corregido}} = \text{FactorK}_{\text{base}}
  \times \left(\frac{d_{\text{calibración}}}{d_{\text{medición}}}\right)^2$$

\textbf{¿Por qué al cuadrado?} Porque el área es bidimensional.
Si la distancia se duplica, el objeto aparece 2 veces más pequeño en cada dimensión,
por lo tanto ocupa $2^2 = 4$ veces menos píxeles.
\end{formulabox}

\vspace{0.3cm}

\textbf{Ejemplo:} Calibración a 80 cm, medición a 100 cm:
$$\text{FactorK}_{\text{corregido}} = \text{FactorK}_{\text{base}} \times \left(\frac{80}{100}\right)^2 = \text{FactorK}_{\text{base}} \times 0{,}64$$

\subsection{Vigencia y validez de la calibración}

\begin{center}
\begin{tabular}{lll}
\toprule
\textbf{Condición} & \textbf{Umbral} & \textbf{Acción del sistema} \\
\midrule
Antigüedad de la calibración & $>$ 30 días & Advertencia al usuario \\
Diferencia de resolución de cámara & $>$ ±10\% & Advertencia al usuario \\
FactorK fuera de rango & $<10^{-9}$ o $>1000$ & Error — calibración rechazada \\
Área de referencia muy pequeña & $<$ 1000 px² & Error — selección rechazada \\
\bottomrule
\end{tabular}
\end{center}

% ============================================================
\section{Proceso 3: Pipeline de Procesamiento de Imagen}
% ============================================================

Este es el corazón del sistema. El \texttt{OpenCVService} procesa cada fotograma
de la cámara (30 veces por segundo) siguiendo una cadena de transformaciones.

\begin{infobox}[title=¿Qué es OpenCV?]
\textbf{OpenCV} (Open Computer Vision) es una biblioteca de código abierto
que contiene cientos de algoritmos para análisis de imágenes. VisiSize la
usa para convertir colores, detectar bordes y encontrar objetos.
\end{infobox}

\subsection{Flujograma del pipeline de imagen}

\vspace{0.5cm}
\begin{center}
\begin{tikzpicture}[node distance=0.5cm]

  \node[io] (frame) {Frame BGR\\ de la cámara};

  \node[process, below=of frame]
    (gray) {\textbf{Paso 1:} Conversión a escala de grises\\
            $\text{BGR} \rightarrow \text{GRAY}$};

  \node[process, below=of gray]
    (blur) {\textbf{Paso 2:} Desenfoque gaussiano\\
            Kernel $5\times5$, $\sigma=0$};

  \node[process, below=of blur]
    (thresh) {\textbf{Paso 3:} Umbralización adaptativa\\
              Bloque: 11, Constante: 2};

  \node[process, below=of thresh]
    (morph) {\textbf{Paso 4:} Morfología — cierre\\
             Kernel elíptico $5\times5$};

  \node[process, below=of morph]
    (contours) {\textbf{Paso 5:} Detección de contornos externos\\
                \texttt{RETR\_EXTERNAL + CHAIN\_APPROX\_SIMPLE}};

  \node[process, below=of contours]
    (filter) {\textbf{Paso 6:} Filtrado por área\\
              $[0{,}1\%,\ 95\%]$ del área total del frame};

  \node[decision, below=of filter]
    (anyvalid) {¿Hay contornos\\válidos?};

  \node[process, right=2.5cm of anyvalid]
    (noresult) {Retornar resultado\\vacío (sin objeto)};

  \node[process, below=of anyvalid]
    (largest) {\textbf{Paso 7:} Seleccionar el contorno\\de mayor área};

  \node[process, below=of largest]
    (metrics) {\textbf{Paso 8:} Calcular métricas del contorno};

  \node[process, below=of metrics]
    (overlay) {\textbf{Paso 9:} Dibujar overlay visual\\(contorno + bounding box)};

  \node[io, below=of overlay]
    (result) {ProcessingResult\\(métricas + imágenes intermedias)};

  \draw[arrow] (frame)--(gray)--(blur)--(thresh)--(morph)--(contours)--(filter)--(anyvalid);
  \draw[arrow] (anyvalid) -- node[left, font=\scriptsize]{Sí} (largest);
  \draw[arrow] (anyvalid) -- node[above, font=\scriptsize]{No} (noresult);
  \draw[arrow] (largest)--(metrics)--(overlay)--(result);

\end{tikzpicture}
\end{center}

\subsection{Descripción detallada de cada paso}

\subsubsection{Paso 1 — Conversión a escala de grises}

\begin{notebox}[title={¿Por qué escala de grises?}]
Las imágenes a color tienen 3 valores por píxel (Rojo, Verde, Azul).
La escala de grises solo tiene 1 (luminosidad). Procesar en grises es
3 veces más rápido y los algoritmos de detección de bordes funcionan mejor.
\end{notebox}

La operación es:
$$\text{Gris}(x,y) = 0{,}114 \cdot B(x,y) + 0{,}587 \cdot G(x,y) + 0{,}299 \cdot R(x,y)$$

\subsubsection{Paso 2 — Desenfoque gaussiano}

El desenfoque suaviza el ruido de la imagen (granulado de la cámara).
Sin este paso, el detector de bordes encontraría miles de bordes falsos
en el ruido.

Se usa un \textbf{filtro gaussiano} con kernel de $5 \times 5$ píxeles.
La operación convoluciona la imagen con una campana de Gauss:
$$I'(x,y) = \sum_{u=-2}^{2}\sum_{v=-2}^{2} G(u,v) \cdot I(x+u,\ y+v)$$

\subsubsection{Paso 3 — Umbralización adaptativa}

Convierte la imagen de grises a \textbf{blanco y negro puro} (binaria).
¿Negro o blanco? Depende de si el píxel es más oscuro que su vecindario local.

\begin{center}
\begin{tabular}{lp{7cm}}
\toprule
\textbf{Parámetro} & \textbf{Significado} \\
\midrule
Tipo & \texttt{ADAPTIVE\_THRESH\_GAUSSIAN\_C} \\
Inversión & \texttt{THRESH\_BINARY\_INV} — el objeto queda en blanco \\
Tamaño de bloque & 11 — lado del vecindario que define el umbral local \\
Constante C & 2 — se resta a la media para afinar el umbral \\
\bottomrule
\end{tabular}
\end{center}

La umbralización \textbf{adaptativa} (vs. global) es clave porque maneja
bien variaciones de iluminación: sombras, gradientes de luz, reflejos.

\subsubsection{Paso 4 — Morfología: operación de cierre}

Un objeto real puede tener pequeños agujeros blancos dentro (brillo, texturas).
La \textbf{operación de cierre} (dilatación seguida de erosión) rellena esos huecos:

$$\text{Cierre}(I) = \text{Erode}\big(\text{Dilate}(I,\ K),\ K\big)$$

Se usa un kernel elíptico $5 \times 5$, que simula un pincel circular.

\subsubsection{Paso 5 — Detección de contornos externos}

A partir de la imagen binaria limpia, OpenCV traza los \textbf{bordes}
de todas las regiones blancas. Solo se buscan contornos externos
(\texttt{RETR\_EXTERNAL}), ignorando agujeros dentro de objetos.

\subsubsection{Paso 6 — Filtrado por área}

No todos los contornos encontrados son el objeto de interés.
Se descartan los que son demasiado pequeños (ruido) o demasiado grandes
(fondo completo de la imagen):

$$A_{\min} = 0{,}001 \times W \times H \qquad A_{\max} = 0{,}95 \times W \times H$$

donde $W$ y $H$ son el ancho y alto del frame en píxeles.

\subsubsection{Paso 7 — Selección del contorno más grande}

De todos los contornos válidos que quedaron, se selecciona el de mayor área.
Este es el objeto principal que el usuario quiere medir.

\subsubsection{Paso 8 — Cálculo de métricas}

Con el contorno seleccionado se calculan cuatro métricas:

\begin{center}
\begin{tabularx}{\textwidth}{lXl}
\toprule
\textbf{Métrica} & \textbf{Descripción} & \textbf{Fórmula} \\
\midrule
Área (px²) & Número de píxeles que cubre el objeto &
  $A = \texttt{cv2.contourArea(c)}$ \\[0.3cm]
Perímetro (px) & Longitud total del borde &
  $P = \texttt{cv2.arcLength(c)}$ \\[0.3cm]
Circularidad & Qué tan circular es el objeto (1.0 = círculo perfecto) &
  $C = \dfrac{4\pi A}{P^2}$ \\[0.4cm]
Relación de aspecto & Ancho / Alto del bounding box &
  $R = \dfrac{w_{\text{bbox}}}{h_{\text{bbox}}}$ \\
\bottomrule
\end{tabularx}
\end{center}

\subsubsection{Paso 9 — Overlay visual}

El resultado se superpone sobre el frame original en tiempo real:
\begin{itemize}
  \item Contorno del objeto dibujado con color según la categoría A/B/C/D.
  \item Rectángulo delimitador (bounding box).
  \item Texto con el área en m² y la categoría.
\end{itemize}

% ============================================================
\section{Proceso 4: Cálculo del Área Real y Clasificación}
% ============================================================

\subsection{De píxeles a metros cuadrados}

Una vez que el pipeline entrega el área en píxeles, la conversión es directa:

\begin{formulabox}[title=Fórmula central de medición]
$$\text{Área real}\ (\text{m}^2) = \text{Área}\ (\text{px}^2) \times \text{FactorK}_{\text{corregido}}$$

donde:
$$\text{FactorK}_{\text{corregido}} = \text{FactorK}_{\text{base}} \times \left(\frac{d_{\text{cal}}}{d_{\text{med}}}\right)^2$$
\end{formulabox}

\vspace{0.3cm}

\textbf{Ejemplo numérico completo:}
\begin{itemize}
  \item FactorK base calibrado a 80 cm: $1{,}285 \times 10^{-6}$ m²/px²
  \item Distancia de medición actual: 100 cm
  \item Área detectada: 65.000 px²
\end{itemize}

$$\text{FactorK}_{\text{corr}} = 1{,}285\times10^{-6} \times \left(\frac{80}{100}\right)^2 = 8{,}224\times10^{-7}\ \frac{\text{m}^2}{\text{px}^2}$$
$$\text{Área real} = 65000 \times 8{,}224\times10^{-7} \approx 0{,}0535\ \text{m}^2$$

\subsection{Clasificación en zonas de almacenamiento}

El área calculada $A_{\text{m}^2}$ se evalúa primero contra los umbrales
de rango del sistema. Si la pieza está dentro del rango operable, se le
asigna una de las cinco zonas (A--E); en caso contrario se marca como
\textbf{Fuera de Rango} (FdR) \cite{malamas2003, gonzalez2018}.

\subsubsection*{Diagrama de rangos de zona}

\vspace{0.4cm}
\begin{center}
\begin{tikzpicture}[node distance=0.1cm]
  % Fuera de rango inferior
  \fill[catOut] (-2.2,0) rectangle (0,1);
  % Zonas A--E
  \fill[catA]   (0,0)    rectangle (2.5,1);
  \fill[catB]   (2.5,0)  rectangle (5.0,1);
  \fill[catC]   (5.0,0)  rectangle (7.5,1);
  \fill[catD]   (7.5,0)  rectangle (10.0,1);
  \fill[catE]   (10.0,0) rectangle (12.5,1);
  % Fuera de rango superior
  \fill[catOut] (12.5,0) rectangle (14.7,1);

  \node[white, font=\tiny\bfseries]       at (-1.1,  0.5) {FdR inf.};
  \node[white, font=\scriptsize\bfseries] at (1.25,  0.5) {Zona A};
  \node[white, font=\scriptsize\bfseries] at (3.75,  0.5) {Zona B};
  \node[white, font=\scriptsize\bfseries] at (6.25,  0.5) {Zona C};
  \node[white, font=\scriptsize\bfseries] at (8.75,  0.5) {Zona D};
  \node[white, font=\scriptsize\bfseries] at (11.25, 0.5) {Zona E};
  \node[white, font=\tiny\bfseries]       at (13.6,  0.5) {FdR sup.};

  \node[font=\tiny] at (-2.2, -0.5) {$<0{,}0025$};
  \node[font=\tiny] at (-2.2, -0.9) {$\text{m}^2$};
  \node[font=\tiny] at (0,    -0.5) {$0{,}0025\,\text{m}^2$};
  \node[font=\tiny] at (2.5,  -0.5) {$1{,}0$};
  \node[font=\tiny] at (5.0,  -0.5) {$2{,}0$};
  \node[font=\tiny] at (7.5,  -0.5) {$3{,}0$};
  \node[font=\tiny] at (10.0, -0.5) {$4{,}0$};
  \node[font=\tiny] at (10.0, -0.9) {$\text{m}^2$};
  \node[font=\tiny] at (12.5, -0.5) {$\text{FOV}_{\max}$};
\end{tikzpicture}
\end{center}

\vspace{0.5cm}
\begin{center}
\begin{tabular}{p{1.1cm} p{2.5cm} p{3.8cm} p{2.4cm} p{1.8cm}}
\toprule
\textbf{Zona} & \textbf{Nombre} & \textbf{Rango de área} & \textbf{Color} & \textbf{Hex} \\
\midrule
\textbf{FdR} & Fuera de rango &
  $A < 0{,}0025\,\text{m}^2$ \textit{ó} excede FOV &
  \textcolor{catOut}{\rule{1.2cm}{0.25cm}} Rojo & \texttt{\#f44336} \\[3pt]
\textbf{A} & Primera zona &
  $[0{,}0025 - 1{,}0)\,\text{m}^2$ &
  \textcolor{catA}{\rule{1.2cm}{0.25cm}} Verde & \texttt{\#4caf50} \\[3pt]
\textbf{B} & Segunda zona &
  $[1{,}0 - 2{,}0)\,\text{m}^2$ &
  \textcolor{catB}{\rule{1.2cm}{0.25cm}} Azul & \texttt{\#2196f3} \\[3pt]
\textbf{C} & Tercera zona &
  $[2{,}0 - 3{,}0)\,\text{m}^2$ &
  \textcolor{catC}{\rule{1.2cm}{0.25cm}} \'{A}mbar & \texttt{\#ff9800} \\[3pt]
\textbf{D} & Cuarta zona &
  $[3{,}0 - 4{,}0)\,\text{m}^2$ &
  \textcolor{catD}{\rule{1.2cm}{0.25cm}} Naranja & \texttt{\#ff5722} \\[3pt]
\textbf{E} & Quinta zona &
  $\geq 4{,}0\,\text{m}^2$ (cabe en frame) &
  \textcolor{catE}{\rule{1.2cm}{0.25cm}} Morado & \texttt{\#9c27b0} \\
\bottomrule
\end{tabular}
\end{center}

\begin{notebox}[title={Mapeo con el código fuente --- MeasurementCategory}]
En el código Python, la enumeración \texttt{MeasurementCategory} define:
\texttt{A}, \texttt{B}, \texttt{C}, \texttt{D} (equivalentes a las zonas
homónimas) y \texttt{error} (mapea a Zona~E si la pieza cabe en el frame,
o a FdR si excede el campo de visión). Esta distinción se formalizará en
versiones futuras del sistema.
\end{notebox}

\subsection{Advertencia de borde de zona}

Si el área medida está dentro del $\pm5\,\%$ de un límite entre zonas,
el sistema emite una advertencia visual indicando que un pequeño error de
medición podría cambiar la zona de almacenamiento asignada:

$$\text{Advertencia si: } \left|\frac{\text{Área} - L_i}{L_i}\right| \leq 0{,}05$$

donde $L_i \in \{0{,}0025;\; 1{,}0;\; 2{,}0;\; 3{,}0;\; 4{,}0\}$ son los
límites entre zonas (incluido el umbral FdR/Zona~A inferior).

\begin{notebox}[title={Ejemplo de advertencia}]
Si se mide un área de $0{,}972\ \text{m}^2$:
$$\left|\frac{0{,}972 - 1{,}0}{1{,}0}\right| = 0{,}028 < 0{,}05$$
El sistema indicará: \emph{«El retal está cerca del límite Zona~A / Zona~B.
Un pequeño error de medición podría cambiar la zona de almacenamiento asignada.»}
\end{notebox}

% ============================================================
\section{Proceso 5: Gestión de la Cámara}
% ============================================================

\subsection{Flujograma de inicialización de cámara}

\vspace{0.5cm}
\begin{center}
\begin{tikzpicture}[node distance=0.55cm]

  \node[startstop] (s) {Solicitud de inicializar cámara};

  \node[decision, below=0.7cm of s]
    (tipo) {¿URL de\\cámara IP?};

  \node[process, left=2.5cm of tipo]
    (webcam) {Abrir \texttt{cv2.VideoCapture(índice)}\\(webcam local)};

  \node[process, right=2.5cm of tipo]
    (ip) {Validar URL RTSP/HTTP\\Abrir \texttt{cv2.VideoCapture(url)}};

  \node[process, below=1.5cm of tipo]
    (setres) {Configurar resolución:\\$1280 \times 720$ @ 30 fps};

  \node[decision, below=of setres]
    (ok) {¿Cámara\\abierta?};

  \node[io, right=2cm of ok]
    (err) {Lanzar excepción\\con mensaje de error};

  \node[process, below=of ok]
    (testframe) {Capturar frame de prueba};

  \node[decision, below=of testframe]
    (frmok) {¿Frame\\válido?};

  \node[startstop, below=of frmok]
    (listo) {Cámara lista};

  \draw[arrow] (s)--(tipo);
  \draw[arrow] (tipo) -- node[above, font=\scriptsize]{No} (webcam);
  \draw[arrow] (tipo) -- node[above, font=\scriptsize]{Sí} (ip);
  \draw[arrow] (webcam) |- (setres);
  \draw[arrow] (ip) |- (setres);
  \draw[arrow] (setres)--(ok);
  \draw[arrow] (ok) -- node[left, font=\scriptsize]{Sí} (testframe);
  \draw[arrow] (ok) -- node[above, font=\scriptsize]{No} (err);
  \draw[arrow] (testframe)--(frmok);
  \draw[arrow] (frmok) -- node[left, font=\scriptsize]{Sí} (listo);
  \draw[arrow] (frmok.east) -- ++(1.5,0) |- node[right, font=\scriptsize]{No} (err.east);

\end{tikzpicture}
\end{center}

\subsection{Tipos de fuente de imagen}

\begin{center}
\begin{tabularx}{\textwidth}{llX}
\toprule
\textbf{Tipo} & \textbf{Formato} & \textbf{Ejemplo} \\
\midrule
Webcam local & Índice entero & \texttt{0} (cámara integrada), \texttt{1} (USB) \\
Cámara IP — RTSP & URL RTSP & \texttt{rtsp://192.168.1.10:554/stream} \\
Cámara IP — HTTP & URL HTTP & \texttt{http://192.168.1.10:8080/video} \\
\bottomrule
\end{tabularx}
\end{center}

% ============================================================
\section{Proceso 6: Almacenamiento de Datos}
% ============================================================

\subsection{Base de datos SQLite (\texttt{measurements.db})}

SQLite es una base de datos que vive en un único archivo.
No requiere servidor, ni instalación especial. Perfecta para
aplicaciones de escritorio como VisiSize.

\vspace{0.4cm}
\begin{center}
\begin{tikzpicture}
  \node[draw=purple!60, fill=purple!8, rounded corners=6pt,
        minimum width=12cm, font=\small] (tabla) {
    \begin{tabular}{llp{6cm}}
      \toprule
      \textbf{Campo} & \textbf{Tipo} & \textbf{Descripción} \\
      \midrule
      \texttt{id} & INTEGER PK & Identificador autoincrementado \\
      \texttt{name} & TEXT & Nombre opcional del registro \\
      \texttt{area\_pixels} & REAL & Área en píxeles cuadrados \\
      \texttt{area\_m2} & REAL & Área real en m² \\
      \texttt{category} & TEXT & Categoría A/B/C/D/error \\
      \texttt{factor\_k\_used} & REAL & FactorK aplicado \\
      \texttt{distance\_cm} & REAL & Distancia cámara-objeto \\
      \texttt{timestamp} & TEXT & Fecha y hora de la medición \\
      \texttt{notes} & TEXT & Notas libres \\
      \texttt{silhouette\_path} & TEXT & Ruta a imagen PNG \\
      \texttt{perimeter\_pixels} & REAL & Perímetro en píxeles \\
      \texttt{circularity} & REAL & Índice de circularidad \\
      \texttt{confidence} & REAL & Confianza IA (reservado) \\
      \texttt{is\_valid} & INTEGER & 1 = activa, 0 = inválida \\
      \bottomrule
    \end{tabular}
  };
\end{tikzpicture}
\end{center}

\subsection{Guardado de siluetas en \texttt{data/images/}}

Al capturar una medición, la silueta del objeto se guarda como imagen PNG.
Para ahorrar espacio, la imagen se redimensiona si supera 512 px por lado:

\begin{center}
\begin{tikzpicture}[node distance=0.5cm]
  \node[io] (frame) {Frame original con contorno};
  \node[process, below=of frame] (crop) {Recortar región del bounding box};
  \node[decision, below=of crop] (size) {¿Lado mayor\\$>$ 512 px?};
  \node[process, right=2cm of size] (resize) {Redimensionar\\(mantener proporción)};
  \node[process, below=of size] (save) {Guardar como PNG\\(compresión nivel 6)};
  \node[storage, below=of save] (disk) {\texttt{data/images/}};

  \draw[arrow] (frame)--(crop)--(size);
  \draw[arrow] (size) -- node[above, font=\scriptsize]{Sí} (resize);
  \draw[arrow] (resize) |- (save);
  \draw[arrow] (size) -- node[left, font=\scriptsize]{No} (save);
  \draw[arrow] (save)--(disk);
\end{tikzpicture}
\end{center}

% ============================================================
\section{Proceso 7: Interfaz de Usuario}
% ============================================================

La UI está construida con \textbf{CustomTkinter}, una versión mejorada de la
biblioteca estándar de Python para interfaces gráficas, que ofrece un aspecto
moderno con soporte para tema oscuro y claro.

\subsection{Estructura de la interfaz}

\begin{center}
\scalebox{0.88}{%
\begin{tikzpicture}[node distance=0.4cm]
  \node[draw=visiblu, fill=visiblu!10, rounded corners=6pt,
        minimum width=12cm, minimum height=0.7cm, font=\small\bfseries]
        (mw) {MainWindow — Ventana principal};

  \node[draw=blue!50, fill=blue!6, rounded corners=4pt,
        minimum width=3.6cm, minimum height=4.5cm, font=\small,
        below left=1.2cm and 4.0cm of mw, align=center, text width=3.2cm]
        (cam) {\textbf{Pestaña «Cámara»}\\[0.2cm]
               Feed en tiempo real\\Botón iniciar/detener\\Control de distancia\\Botón «Capturar»\\Resultado actual};

  \node[draw=green!50!black, fill=green!6, rounded corners=4pt,
        minimum width=3.6cm, minimum height=4.5cm, font=\small,
        below=1.2cm of mw, align=center, text width=3.2cm]
        (hist) {\textbf{Pestaña «Historial»}\\[0.2cm]
                Lista de mediciones\\Detalles por registro\\Imagen de silueta\\Eliminar individual\\Borrar todo};

  \node[draw=purple!50, fill=purple!6, rounded corners=4pt,
        minimum width=3.6cm, minimum height=4.5cm, font=\small,
        below right=1.2cm and 4.0cm of mw, align=center, text width=3.2cm]
        (cfg) {\textbf{Pestaña «Configuración»}\\[0.2cm]
               Fuente de cámara\\Gestión de calibración\\Ajuste de tema\\Versión de app};

  \draw[arrow] (mw.south) -- (cam.north);
  \draw[arrow] (mw.south) -- (hist.north);
  \draw[arrow] (mw.south) -- (cfg.north);
\end{tikzpicture}}
\end{center}

\subsection{Flujo de interacción en la pestaña «Cámara»}

\vspace{0.4cm}
\begin{center}
\begin{tikzpicture}[node distance=0.5cm]

  \node[io] (user) {Usuario abre pestaña «Cámara»};
  \node[process, below=of user] (btn) {Presiona «Iniciar cámara»};
  \node[process, below=of btn] (stream) {Hilo de captura activo:\\30 frames/segundo};
  \node[process, below=of stream] (disp) {Cada frame procesado y\\mostrado en el widget};
  \node[decision, below=of disp] (adj) {¿Usuario ajusta\\distancia?};
  \node[process, right=2.5cm of adj] (recalc) {Recalcular FactorK\\corregido};
  \node[decision, below=of adj] (cap) {¿Usuario presiona\\«Capturar»?};
  \node[process, below=of cap] (dialog) {Diálogo de confirmación\\(nombre + notas opcionales)};
  \node[process, below=of dialog] (store) {Guardar en SQLite + PNG};
  \node[io, below=of store] (ok) {Notificación de éxito\\en la UI};

  \draw[arrow] (user)--(btn)--(stream)--(disp)--(adj);
  \draw[arrow] (adj) -- node[above, font=\scriptsize]{Sí} (recalc);
  \draw[arrow] (recalc) |- (disp);
  \draw[arrow] (adj) -- node[left, font=\scriptsize]{No} (cap);
  \draw[arrow] (cap) -- node[left, font=\scriptsize]{Sí} (dialog);
  \draw[arrow] (dialog)--(store)--(ok);
  \draw[warrow] (cap.east) -- ++(1.2,0) |- node[right, font=\scriptsize]{No} (disp.east);

\end{tikzpicture}
\end{center}

% ============================================================
\section{Proceso 8: Módulo IA — Clasificador YOLO (En desarrollo)}
% ============================================================

VisiSize incluye la infraestructura para un clasificador basado en
\textbf{YOLOv8} (You Only Look Once, versión 8), un modelo de visión
artificial de vanguardia.

\begin{infobox}[title=¿Qué es YOLO?]
\textbf{YOLO} es un modelo de redes neuronales convolucionales que puede
identificar y clasificar objetos en una imagen en tiempo real.
En el contexto de VisiSize, identificaría el \emph{tipo} de objeto
(lámina, tubo, panel, placa...) no solo su área.
\end{infobox}

\subsection{Estado actual}

\begin{center}
\begin{tabular}{ll}
\toprule
\textbf{Componente} & \textbf{Estado} \\
\midrule
\texttt{classifier\_interface.py} & Interfaz abstracta definida ✓ \\
\texttt{yolo\_classifier.py} & Esqueleto implementado ✓ \\
Modelo entrenado (\texttt{.pt}) & Pendiente — en entrenamiento \\
Integración con medición & Pendiente — se activará sin cambios en la UI \\
\bottomrule
\end{tabular}
\end{center}

\subsection{Flujo futuro del clasificador}

\vspace{0.4cm}
\begin{center}
\begin{tikzpicture}[node distance=0.5cm]
  \node[io] (img) {Imagen del objeto detectado};
  \node[process, below=of img] (yuv) {Pre-procesar imagen\\para YOLOv8};
  \node[process, below=of yuv] (infer) {Inferencia del modelo YOLOv8\\(\texttt{model.predict()})};
  \node[process, below=of infer] (parse) {Parsear resultados:\\clase, confianza, bounding box};
  \node[io, below=of parse] (out) {Tipo de objeto + Confianza (\%)};

  \draw[arrow] (img)--(yuv)--(infer)--(parse)--(out);
\end{tikzpicture}
\end{center}

% ============================================================
\section{Flujo Completo de una Medición (Resumen Integrado)}
% ============================================================

Esta sección integra todos los procesos en un único flujograma de alto nivel.

\vspace{0.5cm}
\begin{center}
\begin{tikzpicture}[node distance=0.45cm and 1.8cm]

  \node[startstop] (start) {Usuario inicia la cámara};

  \node[process, below=of start]
    (frame) {Capturar frame BGR};

  \node[process, below=of frame]
    (proc) {Pipeline OpenCV:\\grises → blur → umbral → morfología\\→ contornos → filtrado};

  \node[decision, below=of proc]
    (det) {¿Objeto\\detectado?};

  \node[io, right=2cm of det]
    (nodet) {Overlay «sin objeto»\\en la pantalla};

  \node[process, below=of det]
    (metrics) {Calcular métricas:\\área px², perímetro, circularidad};

  \node[process, below=of metrics]
    (dist) {Leer distancia actual\\del panel de control};

  \node[process, below=of dist]
    (fk) {Aplicar corrección de distancia\\al FactorK};

  \node[process, below=of fk]
    (area) {$\text{Área m}^2 = A_{px} \times \text{FactorK}_{\text{corr}}$};

  \node[decision, diamond, aspect=2, below=0.6cm of area, font=\small]
    (fdr) {\textbf{?Fuera de}\\\ \textbf{rango?}};
  \node[io, right=2.0cm of fdr, text width=2.6cm, align=center,
        font=\small]
    (fdrout) {Alerta roja:\\ pieza FdR};
  \node[process, below=0.6cm of fdr]
    (cat) {Asignar zona: A\,/\,B\,/\,C\,/\,D\,/\,E};

  \node[decision, below=of cat]
    (bound) {¿Cerca de\\límite ±5\%?};

  \node[io, right=2cm of bound]
    (warn) {Mostrar advertencia\\de borde de zona};

  \node[process, below=of bound]
    (show) {Mostrar overlay con resultado\\en tiempo real};

  \node[decision, below=of show]
    (capbtn) {¿Botón\\«Capturar»?};

  \node[process, below=of capbtn]
    (confirm) {Diálogo de confirmación};

  \node[process, below=of confirm]
    (dbsave) {INSERT en measurements.db};

  \node[process, below=of dbsave]
    (pngsave) {Guardar silueta PNG};

  \node[process, below=of pngsave]
    (log) {Registrar en app.log};

  \node[startstop, below=of log]
    (done) {Medición guardada exitosamente};

  % Flechas
  \draw[arrow] (start)--(frame)--(proc)--(det);
  \draw[arrow] (det) -- node[left,font=\scriptsize]{Sí} (metrics);
  \draw[arrow] (det) -- node[above,font=\scriptsize]{No} (nodet);
  \draw[arrow] (nodet) |- (frame);
  \draw[arrow] (metrics)--(dist)--(fk)--(area)--(fdr);
  \draw[arrow] (fdr) -- node[right, font=\scriptsize]{No} (cat);
  \draw[arrow] (fdr) -- node[above, font=\scriptsize]{S\'{i}} (fdrout);
  \draw[arrow] (fdrout) |- (frame);
  \draw[arrow] (cat) -- (bound);
  \draw[arrow] (bound) -- node[left,font=\scriptsize]{No} (show);
  \draw[arrow] (bound) -- node[above,font=\scriptsize]{Sí} (warn);
  \draw[arrow] (warn) |- (show);
  \draw[arrow] (show)--(capbtn);
  \draw[arrow] (capbtn) -- node[left,font=\scriptsize]{Sí} (confirm);
  \draw[arrow] (confirm)--(dbsave)--(pngsave)--(log)--(done);
  \draw[warrow] (capbtn.east) -- ++(1.5,0) |- node[right,font=\scriptsize]{No} (frame.east);

\end{tikzpicture}
\end{center}

% ============================================================
\section{Estructura de Archivos del Proyecto}
% ============================================================

\begin{lstlisting}[style=pycode, language=bash, caption=Estructura de directorios del proyecto]
AreaCamPython_v1/
|
+-- config.py              # Configuracion global (camara, rutas, tema)
+-- main.py                # Punto de entrada de la aplicacion
+-- requirements.txt       # Dependencias Python
+-- app.log                # Log de actividad (se crea al ejecutar)
|
+-- core/                  # Nucleo puro: sin I/O ni dependencias de UI
|   +-- constants.py       # Parametros del pipeline OpenCV y limites de categoria
|   +-- entities.py        # Estructuras de datos (Medicion, Calibracion, categorias)
|   +-- formulas.py        # Calculos puros (FactorK, area, categorias, correccion)
|   +-- exceptions.py      # Excepciones propias del dominio
|
+-- services/              # Logica de negocio y orquestacion
|   +-- calibration_service.py  # Gestion del FactorK y archivo JSON
|   +-- measurement_service.py  # Orquesta el procesamiento de cada frame
|   +-- camera_service.py       # Conexion con webcam o camara IP
|   +-- opencv_service.py       # Pipeline OpenCV completo
|   +-- ai_classifier/
|       +-- classifier_interface.py  # Contrato abstracto del clasificador
|       +-- yolo_classifier.py       # Implementacion YOLO (en desarrollo)
|
+-- data/
|   +-- calibration.json   # Calibracion activa
|   +-- measurements.db    # Historial de mediciones (SQLite)
|   +-- images/            # Siluetas PNG guardadas
|
+-- ui/
    +-- app.py             # Bootstrap de la aplicacion
    +-- theme.py           # Paleta de colores modo oscuro/claro
    +-- windows/           # Ventanas principales
    |   +-- main_window.py
    |   +-- calibration_window.py
    |   +-- measurement_detail_window.py
    +-- widgets/           # Componentes reutilizables
        +-- camera_feed_widget.py
        +-- area_selector_widget.py
        +-- measurement_card.py
\end{lstlisting}

% ============================================================
\section{Resumen de Fórmulas Matemáticas}
% ============================================================

\begin{center}
\renewcommand{\arraystretch}{2.2}
\begin{longtable}{p{4cm} p{5.5cm} p{4cm}}
\toprule
\textbf{Fórmula} & \textbf{Expresión matemática} & \textbf{Uso} \\
\midrule
FactorK base &
  $FK = \dfrac{A_{\text{real}}}{A_{\text{px}}}$ &
  Calibración \\
\midrule
Corrección por distancia &
  $FK' = FK \times \left(\dfrac{d_{\text{cal}}}{d_{\text{med}}}\right)^2$ &
  Ajuste en medición \\
\midrule
Área real &
  $A_{\text{m}^2} = A_{\text{px}} \times FK'$ &
  Conversión central \\
\midrule
Circularidad &
  $C = \dfrac{4\pi A}{P^2}$ &
  Forma del objeto \\
\midrule
Rango mínimo de contorno &
  $A_{\min} = 0{,}001 \times W \times H$ &
  Filtrado de ruido \\
\midrule
Rango máximo de contorno &
  $A_{\max} = 0{,}95 \times W \times H$ &
  Filtrado de fondo \\
\midrule
Umbral de borde de categoría &
  $\left|\dfrac{A - L_i}{L_i}\right| \leq 0{,}05$ &
  Advertencia al usuario \\
\bottomrule
\end{longtable}
\end{center}

% ============================================================
\section{Glosario}
% ============================================================

\begin{description}
  \item[\textbf{FactorK}] Factor de conversión entre píxeles cuadrados y metros cuadrados.
        Es el corazón de la calibración y de todas las mediciones.
  \item[\textbf{Frame}] Un fotograma individual de video. A 30 fps (fotogramas por segundo),
        el programa procesa 30 imágenes por segundo.
  \item[\textbf{Bounding box}] El rectángulo más pequeño que contiene completamente al objeto
        detectado. Se usa para calcular la relación de aspecto y posición.
  \item[\textbf{Contorno}] La curva cerrada que traza el borde exterior de un objeto en una
        imagen binaria.
  \item[\textbf{Umbralización adaptativa}] Proceso de convertir una imagen en escala de grises
        a blanco y negro, donde el umbral varía según el vecindario local de cada píxel.
  \item[\textbf{Morfología (cierre)}] Operación matemática sobre imágenes binarias que rellena
        pequeños huecos interiores sin cambiar el contorno externo del objeto.
  \item[\textbf{SQLite}] Motor de base de datos ligero que guarda todos los datos en un único
        archivo \texttt{.db}, sin necesidad de servidor.
  \item[\textbf{YOLO}] \emph{You Only Look Once} — arquitectura de red neuronal para detección
        de objetos en tiempo real. Versión 8 (YOLOv8) es la utilizada en este proyecto.
  \item[\textbf{Pipeline}] Cadena de procesamiento donde la salida de cada paso
        es la entrada del siguiente.
  \item[\textbf{Overlay}] Información gráfica superpuesta sobre el video en tiempo real
        (contornos, texto, colores de categoría).
  \item[\textbf{OpenCV}] Open Computer Vision Library — biblioteca de visión por computador
        de código abierto, usada para todas las operaciones de imagen.
  \item[\textbf{RTSP}] \emph{Real-Time Streaming Protocol} — protocolo para transmisión de video
        en red, utilizado por cámaras IP.
  \item[\textbf{Retal metálico}] Fragmento sobrante de metal que resulta
  de procesos de corte o troquelado en plantas como ServiMetales de Colombia S.A.S.
  Estos retales son el objeto de medición del sistema VisiSize.

  \item[\textbf{Zona de almacenamiento}] Categoría de destino físico
  asignada a cada retal según su área. El sistema define cinco zonas (A--E)
  diferenciadas por color digital, que se integrarán a futuro en el gemelo digital
  del layout de planta.

  \item[\textbf{Gemelo digital (Digital Twin)}] Representación virtual y
  dinámica del layout físico de la planta, sincronizada en tiempo real con
  los datos de producción~\cite{tao2018}. Las zonas A--E del sistema VisiSize
  corresponden a ubicaciones planificadas dentro de este gemelo.

  \item[\textbf{FOV (Field of View / Campo visual)}] Área máxima visible
  por la cámara desde la distancia de instalación configurada. Define el límite
  superior de medición: objetos que excedan el 95\,\% del FOV se clasifican
  como Fuera de Rango.

  \item[\textbf{FdR (Fuera de Rango)}] Condición que se activa cuando el área
  calculada es menor que 0{,}0025\,m\textsuperscript{2} (5\,cm\,$\times$\,5\,cm)
  o cuando el objeto supera el 95\,\% del campo visual de la cámara. En ambos
  casos el sistema emite una alerta visual en rojo y no asigna zona.

\end{description}

% ============================================================
\section{Referencias Técnicas}

\renewcommand{\refname}{}
\begin{thebibliography}{14}

\bibitem{bradski2008}
G. Bradski and A. Kaehler, \textit{Learning OpenCV: Computer Vision with the OpenCV Library}.
Sebastopol, CA: O'Reilly Media, 2008.

\bibitem{gonzalez2018}
R. C. Gonzalez and R. E. Woods, \textit{Digital Image Processing}, 4th\,ed.
New York, NY: Pearson, 2018.

\bibitem{malamas2003}
E. N. Malamas, E. G. M. Petrakis, M. Zervakis, L. Petit, and J.-D. Legat,
``A survey on industrial machine vision systems,''
\textit{Image Vis. Comput.}, vol.~21, no.~2, pp.~171--188, Feb. 2003.

\bibitem{chin1982}
R. T. Chin and C. A. Harlow,
``Automated visual inspection: A survey,''
\textit{IEEE Trans. Pattern Anal. Mach. Intell.}, vol.~PAMI-4, no.~6,
pp.~557--573, Nov. 1982.

\bibitem{suzuki1985}
S. Suzuki and K. Abe,
``Topological structural analysis of digitized binary images by border following,''
\textit{Comput. Vis. Graph. Image Process.}, vol.~30, no.~1, pp.~32--46, 1985.

\bibitem{bradley2007}
D. Bradley and G. Roth,
``Adapting thresholding using the integral image,''
\textit{J. Graph. Tools}, vol.~12, no.~2, pp.~13--21, 2007.

\bibitem{haralick1992}
R. M. Haralick and L. G. Shapiro, \textit{Computer and Robot Vision}, vol.~1.
Reading, MA: Addison-Wesley, 1992.

\bibitem{redmon2016}
J. Redmon, S. Divvala, R. Girshick, and A. Farhadi,
``You only look once: Unified, real-time object detection,''
in \textit{Proc. IEEE Conf. Comput. Vis. Pattern Recognit. (CVPR)},
Las Vegas, NV, Jun. 2016, pp.~779--788.

\bibitem{jocher2023}
G. Jocher, A. Chaurasia, and J. Qiu,
``Ultralytics YOLO,'' version 8.0, 2023. [Online].
Available: \url{https://github.com/ultralytics/ultralytics}

\bibitem{tao2018}
F. Tao, J. Cheng, Q. Qi, M. Zhang, H. Zhang, and F. Sui,
``Digital twin-driven product design, manufacturing and service with big data,''
\textit{Int. J. Adv. Manuf. Technol.}, vol.~94, no.~9--12,
pp.~3563--3576, Feb. 2018.

\bibitem{li2020}
Y. Li, H. Zhao, S. Hu, Q. Wang, and H. Chen,
``Metal surface defect detection based on improved YOLOv3 neural network,''
\textit{IEEE Access}, vol.~8, pp.~71864--71875, 2020.

\bibitem{owens2010}
M. Owens and G. Allen, \textit{The Definitive Guide to SQLite}, 2nd\,ed.
New York, NY: Apress, 2010.

\bibitem{opencv_docs}
OpenCV Development Team, ``OpenCV Documentation,'' 2024. [Online].
Available: \url{https://docs.opencv.org/}

\bibitem{schimansky2023}
T. Schimansky, ``CustomTkinter Documentation,'' 2023. [Online].
Available: \url{https://customtkinter.tomschimansky.com/}

\end{thebibliography}

\end{document}
